\title{Svensk pølseret}
\PrepTime{20}
\CookingTime{}
\CookingTempe{}
\TypeCooking{Gris}
\NbPerson{4}
\Image{0 100 800 900}{images/svensk_pølseret.jpg}

\begin{ingredient}
    \begin{main}
        \item 1dl Fløde\index[ingredient]{Fløde}
        \item 800g Kogte Kartofler\index[ingredient]{Kartofler}
        \item 1 Løg\index[ingredient]{Løg}\index[hovedret]{Løg}
        \item 2dl Mælk\index[ingredient]{Mælk}
        \item 3 tsk. Paprika\index[ingredient]{Paprika}
        \item Peber\index[ingredient]{Peber}
        \item Persille\index[ingredient]{Persille}
        \item Purløg\index[ingredient]{Purløg}
        \item 500g Pølser\index[ingredient]{Pølser}
        \item Salt\index[ingredient]{Salt}
        \item Smør\index[ingredient]{Smør}
        \item 3 spsk. Tomatpuré\index[ingredient]{Tomatpuré}
    \end{main}
\end{ingredient}
\begin{recipe}
    \step{Pil og hak løget fint.}
    \step{Skær de kogte kartofler ud i grove, mundrette stykker.}
    \step{Skær pølserne i skiver på ca. 1cm.}
    \step{Varm en gryde op ved middel varme med lidt smør eller olie og steg løgene klare heri.}
    \step{Kom pølser og paprika i gryden og svits til pølserne har fået stegeskorpe. Tilsæt herefter tomatpuré, kartofler, mælk og fløde.}
    \step{Lad retten simre ved lavt blus til kartoflerne begynder at koge let ud --- ca. 10 minutter.}
    \step{Smag til med salt og peber.}
    \step{Server med frisk hakket purløg eller persille.}
\end{recipe}
\begin{notes}

\end{notes}

\setcounter{step}{0}
\setcounter{substep}{0}

\clearpage